\documentclass{article}

% Allow long file paths / job IDs in \url and \nolinkurl to break at hyphens.
\PassOptionsToPackage{hyphens}{url}

% NeurIPS 2025 style — swap for iclr2026 when targeting ICLR
\usepackage[final,main]{neurips_2025}

\usepackage[utf8]{inputenc}
\usepackage[T1]{fontenc}
\usepackage{hyperref}
\usepackage{url}
\usepackage{booktabs}
\usepackage{amsfonts}
\usepackage{amsmath}
\usepackage{amssymb}
\usepackage{nicefrac}
\usepackage{microtype}
\usepackage{xcolor}
\usepackage{graphicx}
\usepackage{subcaption}
\usepackage{algorithm}
\usepackage{algpseudocode}
\usepackage{multirow}

% ── convenience macros ────────────────────────────────────────────────────────
\newcommand{\fij}{f_{i,j}}
\newcommand{\mlpedge}{\textsc{MLPEdge}}
\newcommand{\kanprey}{\textsc{KANprey}}
\newcommand{\katprey}{\textsc{KATprey}}
\newcommand{\guppylm}{\textsc{GuppyLM}}
\newcommand{\kanlinear}{\textsc{KANLinear}}
\newcommand{\grkan}{\textsc{GR-KAN}}
\newcommand{\basiskan}{\textsc{BasisKAN}}

% ──────────────────────────────────────────────────────────────────────────────

\title{KAN for Language Models: Empirical Interpretability and the MLPEdge Variant}

\author{%
  Felippe Alves \\
  Telus Digital Research Hub -- CIAAM \\
  Institute of Mathematics and Statistics \\
  University of São Paulo \\
  \texttt{felippe.pereira@usp.br}
  \And
  Renato Vicente \\
  Telus Digital Research Hub -- CIAAM \\
  Institute of Mathematics and Statistics \\
  University of São Paulo \\
  \texttt{rvicente@usp.br}
}

\begin{document}

\maketitle

% ── Abstract ──────────────────────────────────────────────────────────────────
\begin{abstract}
Kolmogorov-Arnold Networks (KANs) replace fixed activation functions with
learnable univariate functions on every edge, yielding~$\sim$147\,K visualizable
1D curves per language model layer and a strong interpretability promise.
We present the first empirical study of KAN edge interpretability at language
model scale, measuring nonlinearity, activity, roughness, and functional
principal components across all edges of KAN-based transformers trained on a
60\,K-sample instruction-following corpus.
Our study finds that \textbf{87.8\%} of edges are meaningfully nonlinear
(NLS~$>0.1$) and only 0.4\% are dead, providing the first data-driven
assessment of the KAN interpretability claim in a language model.
A Mann-Whitney U test on all 884,736 edges yields $p < 10^{-300}$ for KAN\,>\,MLP,
though this $p$-value should be read descriptively rather than inferentially: edges
within the same layer share the same B-spline grid and basis, violating the
independence assumption and inflating the effective sample size.
The quantitative gap is modest (median NLS 0.281 vs.\ 0.268), and the per-layer
distributions are consistent, suggesting the directional result is robust.

Motivated by the observation that the KAN topology---additive decomposition
over edges---may matter more than the specific B-spline basis within the
KAN-family design space, we introduce \mlpedge{}: a drop-in replacement for the
spline edge function with a tiny per-edge MLP ($\mathbb{R} \to \mathbb{R}$).
We also run a follow-up function-basis screen over rational, polynomial,
Gaussian/RBF, compact-support, and ReLU-power edge bases.
On \guppylm{}, the strongest new basis row is a grouped Chebyshev degree-3 FFN,
which reaches $0.2770_{\pm 0.0003}$ validation loss over three seeds,
statistically indistinguishable from the same-protocol MLP baseline
($0.2770_{\pm 0.0008}$) but slower in the same-device inference sweep
(550\,ms vs.\ 309\,ms).
The corrected canonical \grkan{} row remains close to MLP
($0.2780_{\pm 0.0005}$), and a square-denominator rational ablation is similar
($0.2780_{\pm 0.0009}$).
\mlpedge{} remains the strongest \emph{edge-function} replacement among the
B-spline and \mlpedge{} rows ($0.2854_{\pm 0.0005}$ vs.\
$0.2883_{\pm 0.0005}$ for the standard-attention B-spline baseline), while the
Gaussian/RBF and compact-support bases fail the local quality screen.
Thus no tested KAN-family FFN dominates the standard MLP FFN on both quality and
latency at this scale.
A GPT-2 scale experiment on Wikitext-103 (124M parameters, 20\,000 steps)
confirms the scale caveat: the MLP baseline reaches perplexity 16.58 versus
20.96 for the parameter-matched \mlpedge{} ($h=8$), widening rather than
reversing the GuppyLM gap.
For \grkan{}, the analogous scale test is also negative. Training from scratch at
identical scale (286M parameters, ClimbMix corpus, H100 80GB HBM3) under the
canonical Safe Pad\'{e} rational definition, we sweep the group count: the best
corrected configuration, $g=4$, reaches val\_bpb = 0.890 versus 0.8464 for the MLP
baseline (gap $+0.044$, $\sim5.2\%$), while $g=8$ is worse
($0.9030$, $+0.057$, $\sim6.7\%$)---a higher group count gives no benefit at this scale.
On DCLM CORE (22 ICL tasks, centered accuracy), $g=4$ scores 0.128 and $g=8$ scores
0.116, both below the MLP baseline of 0.154 ($\Delta = -0.026$ and $-0.038$); MLP
leads on 18 of 22 tasks even against the stronger $g=4$ variant.
These results are notably weaker than preliminary estimates reported in an earlier
version of this work, which were inflated by an implementation error in the rational
activation denominator (per-coefficient absolute value instead of the canonical
absolute-value-after-sum).
We document this correction transparently.
The corrected gap is wider than the $\sim1.2\%$ preliminary estimate, so we
treat this as a cautionary scale result: rational activations remain less
damaging than the per-edge \mlpedge{} factorization at this scale, but the current
evidence does not support a \grkan{} scaling advantage over MLP at any tested group
count.
All code and checkpoints are released at
\url{https://github.com/HCAI-USP/kan-guppylm}.
\end{abstract}

% ── 1. Introduction ───────────────────────────────────────────────────────────
\section{Introduction}
\label{sec:intro}

Kolmogorov-Arnold Networks~\citep{liu2024kan} replace the fixed activation
functions of standard MLPs with learnable univariate functions on each edge,
grounded in the Kolmogorov-Arnold representation theorem.
The appeal for interpretability is direct: every edge $\fij(x)$ is a
visualizable 1D function that can, in principle, be symbolically regressed
into a closed-form expression.
Subsequent work adapts KANs to transformers~\citep{yang2024kat} and improves
their computational efficiency~\citep{blealtan2024efficientkan}, but the
interpretability promise remains largely theoretical.

\textbf{Two empirical gaps.}
First, no study has measured what KAN edge functions actually learn in a
language model: are they nonlinear, active, and structurally diverse, or do
most edges collapse to near-linear functions?
Second, no work has asked whether the B-spline basis is necessary for quality,
or whether the KAN topology alone is responsible for the gains over MLP
transformers.

\textbf{This paper.}
We address both gaps using \guppylm{}~\citep{guppylm60k} as a controlled
testbed---a 60\,K-sample fish-personality chatbot dataset small enough for
exhaustive edge-level analysis (${\sim}147$\,K functions per FFN layer at full
resolution) yet rich enough to produce coherent chat behavior.
Our contributions are:

\begin{enumerate}
  \item \textbf{Interpretability study.} We reconstruct and statistically
    characterize all learned 1D edge functions in KAN-based transformers,
    reporting nonlinearity scores (NLS), activity, roughness, functional PCA
    archetypes, and a direct comparison with MLP effective edge functions via
    Mann-Whitney~U test (Section~\ref{sec:interpretability}).

  \item \textbf{\mlpedge{}.} We propose replacing B-spline edge functions with
    tiny per-edge MLPs ($\mathbb{R}\to\mathbb{R}$), preserving the KAN
    topology while eliminating the grid mechanism.
    Within the B-spline/\mlpedge{} subset, \mlpedge{} improves over the
    B-spline baselines on validation loss and heuristic prompt quality; a
    same-protocol MLP baseline and the new Chebyshev basis screen are stronger
    on the core training objective (Section~\ref{sec:results}).

  \item \textbf{Function-basis screen.} We test rational denominator
    variants, Chebyshev/Legendre polynomial bases, Gaussian/RBF bases,
    compact-support bases, and ReLU-power bases.
    A Chebyshev degree-3 grouped basis matches the MLP validation-loss mean but
    is slower; Gaussian/RBF and hat-style localized bases do not pass the
    GuppyLM quality screen (Section~\ref{sec:basis_screen}).
\end{enumerate}

% ── 2. Background ─────────────────────────────────────────────────────────────
\section{Background}
\label{sec:background}

\subsection{KAN and the Kolmogorov-Arnold Theorem}

The Kolmogorov-Arnold representation theorem states that every continuous
multivariate function $f: [0,1]^n \to \mathbb{R}$ can be written as:
\begin{equation}
  f(\mathbf{x}) = \sum_{q=0}^{2n} \Phi_q\!\left(\sum_{p=1}^{n}
    \phi_{q,p}(x_p)\right)
\end{equation}
where $\phi_{q,p}$ and $\Phi_q$ are continuous univariate functions.
\citet{liu2024kan} use this as motivation for replacing MLP weight matrices
with learned function arrays, noting that the exact theorem applies to
two-layer networks but the architecture generalizes to arbitrary depth.

\subsection{KANLinear (EfficientKAN)}
\label{sec:kanlinear}

\citet{blealtan2024efficientkan} implement KAN layers efficiently using
B-spline basis functions:
\begin{align}
  \text{forward}(\mathbf{x}) &= \mathbf{x} \odot \text{SiLU}(\mathbf{x})\,
    \mathbf{W}_\text{base}^\top
    + \underbrace{\text{einsum}(\text{"bik,oik}\to\text{bo"},\,
    B(\mathbf{x}),\,
    \mathbf{W}_\text{spline})}_{\text{spline path}}
\end{align}
where $B(\mathbf{x}) \in \mathbb{R}^{B \times d_\text{in} \times n_\text{basis}}$
is the B-spline basis evaluated via the Cox-de~Boor recursion,
$\mathbf{W}_\text{spline} \in \mathbb{R}^{d_\text{out} \times d_\text{in} \times
n_\text{basis}}$, and $n_\text{basis} = \text{grid\_size} + \text{spline\_order}$.
Grid knots adapt to the activation distribution after a warm-up phase
(\emph{grid update}).

\subsection{KAT Attention}
\label{sec:kat}

\citet{yang2024kat} replace the linear Q and K projections with KANLinear
feature maps $\varphi$, turning the dot-product similarity into a learned
kernel:
\begin{equation}
  K(\mathbf{q}, \mathbf{k}) = \varphi(\mathbf{q}) \cdot \varphi(\mathbf{k})
\end{equation}

% ── 3. Models ─────────────────────────────────────────────────────────────────
\section{Models}
\label{sec:models}

All models share the same base transformer: 6 layers, $d_\text{model}=384$,
6 attention heads, trained on the \guppylm{} corpus with prompt-response
masking (loss on assistant tokens only).

\paragraph{MLP-4x-GELU baseline.}
For the controlled small-scale comparison in Section~\ref{sec:results}, the
standard transformer FFN is \texttt{Linear(384$\to$1536)} $\to$ \texttt{GELU}
$\to$ \texttt{Linear(1536$\to$384)}.

\paragraph{\kanprey{}.}
For the controlled small-scale comparison in Section~\ref{sec:results},
\kanprey{} replaces each FFN with a single \kanlinear{} layer
(\texttt{384$\to$384}, grid=2, order=3), producing
$384 \times 384 = 147{,}456$ learnable 1D edge functions per layer
($6 \times 147{,}456 = 884{,}736$ total across all FFN layers).
It uses standard dot-product attention (not KAT).

\paragraph{\katprey{}.}
Extends \kanprey{} by adding KANLinear feature maps to Q and K projections
in every attention head (KAT attention~\citep{yang2024kat}).

\paragraph{\mlpedge{} (proposed).}
Uses standard dot-product attention (same as \kanprey{}, \emph{not} KAT).
Replaces the B-spline edge function $\fij(x)$ with a tiny per-edge MLP:
\begin{align}
  \mathbf{H} &= \sigma\!\left(
    \text{einsum}(\text{"bi,ih}\to\text{bih"},\, \mathbf{x},\, \mathbf{W}_1)
    + \mathbf{b}_1\right)
    && \in \mathbb{R}^{B \times d_\text{in} \times h} \\
  \mathbf{out} &= \text{einsum}(\text{"bih,oih}\to\text{bo"},\,
    \mathbf{H},\, \mathbf{W}_2) + \mathbf{b}_\text{out}
    && \in \mathbb{R}^{B \times d_\text{out}}
\end{align}
where $h$ is the per-edge hidden size (default $h=5$; the B-spline baseline
uses $n_\text{basis}=8$, so $h=8$ is parameter-matched while $h=5$ uses
fewer parameters per edge). The einsum structure is identical to \kanlinear{}, so
$\fij(x)$ is still a visualizable 1D scalar function.
The \texttt{update\_grid()} call is a no-op: no grid mechanism is needed.

\paragraph{\basiskan{} screen.}
For the function-basis experiments, we keep the standard transformer block and
replace the FFN nonlinearity with grouped scalar bases followed by the same
linear mixing dimensions as the MLP FFN.
The screened bases include Chebyshev and Legendre recurrences, Gaussian RBFs,
inverse-quadratic RBFs, Wendland compact-support polynomials, triangular and
quadratic hats, and symmetric ReLU-power functions.
The RBF and hat variants normalize pre-basis activations with \texttt{tanh},
use 8 centers uniformly spaced on $[-1,1]$, and set the Gaussian/RBF width to
\texttt{width\_scale}/7.

Figure~\ref{fig:architecture} illustrates the three primary FFN variants.

% ── 4. Experimental Setup ─────────────────────────────────────────────────────
\section{Experimental Setup}
\label{sec:setup}

\paragraph{Dataset.}
\texttt{arman-bd/guppylm-60k-generic}~\citep{guppylm60k}: 60\,K
instruction-response pairs in ChatML format (fish-personality chatbot).
BPE tokenizer trained on the same data, vocab size 2\,393.

\paragraph{Training.}
Main comparison rows use 8\,000 steps, batch size 32, cosine LR schedule
(peak $3\times10^{-4}$, min $3\times10^{-5}$), AdamW
($\beta_1=0.9$, $\beta_2=0.95$, $\lambda=0.1$), and gradient clipping 1.0.
Prompt-response masking sets targets for the user prompt portion to $-100$
(ignored in cross-entropy loss).
Grid update for KAN/KAT models occurs at step 500.
The function-basis screen uses one seed-42, 2\,000-step pilot for each
candidate; promoted candidates are rerun for 8\,000 steps over seeds
$\{42,43,44\}$.

\paragraph{Evaluation.}
(a)~Validation cross-entropy on a held-out split;
(b)~keyword-match accuracy on 16 fixed prompts with human-defined keyword
sets (pass/fail per prompt);
(c)~mean inference latency over 16 prompts $\times$ 64 new tokens.

\paragraph{Why \guppylm{} as testbed.}
At $d_\text{model}=384$ and 6 layers, each FFN \kanlinear{} layer has
$384\!\times\!384=147{,}456$ edges---small enough to reconstruct and
characterize every 1D function exhaustively, which is infeasible at GPT-2
scale ($d_\text{model}=768$, $4\times$ more edges per layer, $4\times$ more
parameters, orders-of-magnitude more compute per study).
The single-domain corpus provides coherent semantic categories for
token-conditioned interpretability analysis.

% ── 5. Results ────────────────────────────────────────────────────────────────
\section{Results}
\label{sec:results}

\subsection{Main comparison}

\begin{table}[h]
\centering
\scriptsize
\caption{%
  Main \guppylm{} results after protocol cleanup and the function-basis screen.
  Validation loss is mean\,$\pm$\,std over 3 seeds for every 8\,000-step row.
  Eval (16) and latency are measured on seed-42 checkpoints under same-device
  MPS sweeps with 16 prompts and 64 generated tokens; they are heuristic
  behavioral probes rather than replacements for validation loss.
  \katprey{} changes both attention and FFN, so it should not be interpreted
  as an FFN-only ablation.
  The canonical \grkan{} row uses the corrected Safe Pad\'{e} denominator.
}
\label{tab:main}
\begin{tabular}{@{}lcccccc@{}}
\toprule
Model & Attn. & FFN & Val loss & Eval & Lat. & Params \\
\midrule
MLP-4x-GELU & Dot & MLP & $\mathbf{0.2770}_{\pm 0.0008}$ & 11/16 & \textbf{309}\,ms & 11.62M \\
\basiskan{} Cheb $d=3,g=8$ & Dot & Cheb & $\mathbf{0.2770}_{\pm 0.0003}$ & 11/16 & 550\,ms & 11.62M \\
\grkan{} canonical & Dot & Rational & $0.2780_{\pm 0.0005}$ & 14/16 & 459\,ms & 11.62M \\
\grkan{} square denom. & Dot & Rational & $0.2780_{\pm 0.0009}$ & \textbf{15/16} & 382\,ms & 11.62M \\
\kanprey{} grid=2 & Dot & Spline & $0.2883_{\pm 0.0005}$ & 8/16 & 613\,ms & 9.83M \\
\katprey{} & KAT & Spline & $0.2867_{\pm 0.0006}$ & 13/16 & 888\,ms & 10.18M \\
\mlpedge{} $h=5$ & Dot & Edge MLP & $0.2854_{\pm 0.0005}$ & 14/16 & 471\,ms & 8.98M \\
\mlpedge{} $h=8$ & Dot & Edge MLP & $0.2862_{\pm 0.0010}$ & 12/16 & 536\,ms & 11.64M \\
\bottomrule
\end{tabular}
\end{table}

\paragraph{Prompt masking.}
The most impactful single change was adding prompt-response masking
($y_{1:\ell-1} = -100$, where $\ell$ is the assistant prefix length):
computing loss only on assistant tokens dropped \katprey{} validation
loss from 0.4202 to 0.2870 and eliminated role-marker leakage (the model
generating ``\texttt{assistant}'' in responses).
Table~\ref{tab:masking} shows the controlled comparison.

\begin{table}[h]
\centering
\caption{Effect of prompt-response masking on \katprey{} (8\,000 steps).}
\label{tab:masking}
\begin{tabular}{lcc}
\toprule
Training objective & Val loss & Role-marker in responses \\
\midrule
All tokens (no masking) & 0.4202 & Yes \\
Assistant tokens only  & \textbf{0.2870} & No \\
\bottomrule
\end{tabular}
\end{table}

\paragraph{\mlpedge{} relative to B-spline edge-function variants.}
Within the B-spline/\mlpedge{} subset, \mlpedge{} remains the strongest
edge-function replacement we measured.
At $h=5$, it improves mean validation loss over \kanprey{} grid=2
($0.2854$ vs.\ $0.2883$) and over \katprey{} ($0.2867$), while preserving
visualizable 1D edge functions.
This is narrower than a full KAN-family or MLP-replacement claim:
\basiskan{} Cheb $d=3,g=8$ and the MLP baseline are stronger on validation loss,
and MLP remains much faster in the same-device inference sweep.

Table~\ref{tab:ablation} still shows that $h=1$ achieves the lowest
single-seed \mlpedge{} validation loss (0.2828), with loss increasing
monotonically with $h$ on that sweep.
We therefore interpret $h \in \{1,2\}$ as promising \mlpedge{} settings, while
using $h=5$ and $h=8$ for the cleaner controlled comparisons above.

\subsection{Function-basis screen}
\label{sec:basis_screen}

The broader basis screen tests whether the edge-wise KAN topology can recover
MLP-level quality with bases other than B-splines or per-edge MLPs.
Table~\ref{tab:basis_screen} summarizes the seed-42 2\,000-step screen and the
8\,000-step confirmation for promoted candidates.
Only the rational square-denominator ablation and the Chebyshev degree-3 basis
were promoted.

\begin{table}[h]
\centering
\small
\caption{%
  GuppyLM function-basis screen.
  Pilot rows are seed-42, 2\,000-step validation losses.
  Confirmation rows are 8\,000-step mean\,$\pm$\,std over seeds 42/43/44 when
  promoted.
}
\label{tab:basis_screen}
\begin{tabular}{llccc}
\toprule
Family & Config & 2k pilot & 8k confirmation & Decision \\
\midrule
Rational & \grkan{} square denom., $P/Q=5/4,g=8$ & \textbf{0.3134} & $0.2780_{\pm0.0009}$ & confirmed \\
Polynomial & Chebyshev $d=3,g=8$ & 0.3151 & $\mathbf{0.2770}_{\pm0.0003}$ & confirmed \\
Localized & inverse quadratic $c=8,h=1.5$ & 0.4169 & -- & reject \\
Gaussian & $c=8,h=1.0$ & 0.4198 & -- & reject \\
Gaussian & $c=8,h=1.5$ & 0.5123 & -- & reject \\
Gaussian & $c=8,h=2.0$ & 1.1174 & -- & reject \\
Hat/ReLU & ReLU-power $p=2,c=8$ & 0.3664 & -- & reject \\
\bottomrule
\end{tabular}
\end{table}

The main positive result is the Chebyshev row: a recurrence-friendly polynomial
basis matches the MLP validation-loss mean within the observed 3-seed noise,
but it is slower in our unoptimized local implementation (550\,ms vs.\ 309\,ms
for MLP in the seed-42 prompt sweep).
This makes Chebyshev a plausible kernel-engineering target, not an immediate
replacement for the MLP FFN.
The rational square-denominator variant is safe but not a clear improvement over
canonical \grkan{}.

The Gaussian/RBF branch fails for both quality and local kernel reasons.
On shape $4096\times384$ with width scale 1.5, exact \texttt{torch.exp} was
already faster than the tested approximations (2.86\,ms forward / 7.02\,ms
backward) compared with \texttt{exp2} rewrite (3.32\,ms / 8.08\,ms), clamped
exp (3.33\,ms / 8.02\,ms), and LUT interpolation (10.79\,ms / 8.38\,ms).
The LUT was accurate (forward max error $7.7{\times}10^{-6}$; gradient max
error $3.0{\times}10^{-9}$), but slower on MPS.
Training quality was also poor: the best Gaussian width was $h=1.0$ at 0.4198
after 2\,000 steps, far behind Chebyshev and rational variants.


\subsection{Hidden-size ablation}

\begin{table}[h]
\centering
\caption{%
  \mlpedge{} hidden-size ablation (8\,000 steps, seed 42).
  Inference latency measured on Apple M-series (MPS disabled; CPU for
  reproducibility). Bold = best val loss.
}
\label{tab:ablation}
\begin{tabular}{ccc}
\toprule
Hidden $h$ & Val loss & Latency \\
\midrule
\textbf{1}  & \textbf{0.2828} & 226\,ms \\
2  & 0.2836 & 354\,ms \\
5  & 0.2854 & 491\,ms \\
10 & 0.2863 & 593\,ms \\
20 & 0.2886 & 1262\,ms \\
\bottomrule
\end{tabular}
\end{table}

Figure~\ref{fig:ablation} (Appendix) plots the efficiency-quality tradeoff curve.

\subsection{Fair KANprey baseline}
\label{sec:kan_retrain}

The original \kanprey{} result previously highlighted in Table~\ref{tab:main}
(val loss 0.4019) was trained for only 4\,000 steps without prompt-response
masking, making it an unfair comparison point.
We therefore exclude it from the cleaned main table and retain it only as a
protocol-correction example.

Under the corrected 8\,000-step masked setup, the 3-seed \kanprey{} grid=2
baseline reaches $0.2883_{\pm 0.0005}$.
This places it behind \mlpedge{} $h=5$ by 0.0029 validation-loss points and
behind the same-protocol MLP baseline by 0.0113.
The corrected B-spline baseline is still useful as the main controlled
comparison for edge-function parameterizations, but it is no longer a strong
headline model.

\subsection{Scaling to GPT-2 size}
\label{sec:scaling}

To assess whether the small-scale KAN-family improvements of \mlpedge{}
generalize beyond the \guppylm{} testbed, we train a standard MLP transformer
and an \mlpedge{} transformer at GPT-2 small scale
($d_\text{model}=768$, 12 heads, 12 layers) on the Wikitext-103 benchmark
(Table~\ref{tab:scaling}).

Both models were trained for 20\,000 steps with effective batch size 128
(gradient accumulation $4\times32$), cosine LR schedule ($3\times10^{-4}$
to $3\times10^{-5}$, 2\,000 warmup steps), and AdamW with gradient
checkpointing on a single RTX~A6000~48\,GB GPU.
We report the best validation perplexity across all evaluation checkpoints.

\begin{table}[h]
\centering
\caption{%
  GPT-2 scale results on Wikitext-103 (20\,000 steps, effective batch~128).
  Standard MLP FFN uses \texttt{Linear(768$\to$3072)} $\to$ \texttt{GELU}
  $\to$ \texttt{Linear(3072$\to$768)} (4$\times$ expansion).
  \mlpedge{} ($h=8$, matched) is parameter-matched to the MLP baseline,
  isolating the effect of the per-edge factored structure from parameter count.
  Throughput measured end-to-end including gradient accumulation on RTX~A6000.
}
\label{tab:scaling}
\begin{tabular}{lcccc}
\toprule
Model & Val perplexity & Tokens/sec & Train time & Params \\
\midrule
MLP transformer          & \textbf{16.58} & 54{,}820 & 13.3\,h & 124M \\
\mlpedge{} ($h=8$, matched) & 20.96          & 48{,}921 & 14.9\,h & 125M \\
\bottomrule
\end{tabular}
\end{table}

At GPT-2 scale, \mlpedge{} ($h=8$) achieves a validation perplexity of 20.96
versus 16.58 for the MLP baseline---a gap of 4.37 perplexity points (26\%)
in favour of the standard MLP.
This does not reverse an MLP deficit at small scale; rather, it widens a gap
that was already present on \guppylm{}, where MLP was also the strongest
validation-loss baseline after cleanup.
Throughput is 11\% lower for \mlpedge{} due to the \texttt{einsum} overhead
of the per-edge factored computation at $d_\text{model}=768$.

These results suggest that the per-edge factored structure provides an
effective inductive bias at small scale but does not transfer cleanly to
GPT-2 size with the same hidden-per-edge budget.
Whether higher $h$ or a different initialisation scheme would close this gap
is left to future work.
The \grkan{} experiment below tests the related question for a rational-basis variant.

% ── 5.5 GR-KAN at GPT-2 scale ────────────────────────────────────────────────
\subsection{\grkan{} at GPT-2 Scale}
\label{sec:grkan_gpt2}

\paragraph{Setup.}
We train \grkan{} and a standard MLP transformer at the same scale as
Section~\ref{sec:scaling} ($d_\text{model}=768$, 6 heads, 12 layers,
${\sim}286$M parameters) using the nanokan training
harness\footnote{\url{https://github.com/ACS-USP/nanokan}} on the
ClimbMix dataset (a curated web corpus with quality filtering, comparable
to FineWeb-Edu).
Both models are trained for 2\,520 steps (${\sim}1.32$B tokens, effective
batch 256 sequences $\times$ 2\,048 tokens) with the Muon
optimizer~\citep{jordan2024muon} on a single H100 80GB HBM3 GPU.
The \grkan{} FFN uses grouped Pad\'{e} rational activations with polynomial
degrees $P=5$, $Q=4$. We sweep the group count $g \in \{4, 8\}$ (group channel
widths $D/g = 192$ and $96$); a $g=16$ run diverged to NaN and is excluded.
The MLP baseline uses the
standard \texttt{Linear(768$\to$3072)} $\to$ \texttt{SwiGLU}
$\to$ \texttt{Linear(3072$\to$768)} block.
We report validation bits-per-byte (bpb) on a held-out ClimbMix split
and DCLM CORE (22 ICL tasks, centered accuracy), neither of which is
directly comparable to the Wikitext-103 perplexity in
Table~\ref{tab:scaling}, but both isolate the FFN architecture difference under
identical training conditions.

\paragraph{Implementation correction.}
After an earlier version of this work was circulated, a code review identified
an error in the rational activation denominator.
The original Triton kernels in \texttt{rational\_kat\_cu} computed:
\[
Q(x) = 1 + |b_0|\!\cdot\!|x| + |b_1|\!\cdot\!|x|^2 + |b_2|\!\cdot\!|x|^3 + |b_3|\!\cdot\!|x|^4
\]
applying absolute value per coefficient, rather than the canonical Safe
Pad\'{e} rational definition:
\[
D(x) = b_0 x + b_1 x^2 + b_2 x^3 + b_3 x^4,\qquad
Q(x) = 1 + |D(x)|.
\]
This bug was present from the original fork and affected all previously
reported GPT-2 scale \grkan{} results.
The forward and backward Triton kernels have been corrected; gradient
consistency was verified via \texttt{torch.autograd.gradcheck} and a
10-step smoke test on H100 hardware.
The results in Tables~\ref{tab:grkan_gpt2} and~\ref{tab:grkan_core}
reflect the canonical formula.

\paragraph{Engineering: fused Triton backward.}
The na\"{\i}ve PyTorch autograd backward for \texttt{GroupRational} dispatches
${\sim}480$ independent CUDA kernels per training step.
We use a single fused Triton backward kernel covering the Horner
evaluation of $P(x)$, $D(x)$, $D'(x)$, elementwise $\partial_x$ computation, and
block-level coefficient gradient reductions in one launch.
This achieves ${\sim}290$K\,tok/s throughput and ${\sim}86\%$ MFU on H100~80GB~HBM3
(${\sim}3\times$ improvement over the unfused baseline).
The remaining throughput gap to MLP (${\sim}508$K\,tok/s) reflects the arithmetic
intensity of rational activations relative to dense cuBLAS matmuls.

\paragraph{Results.}
Tables~\ref{tab:grkan_gpt2} and~\ref{tab:grkan_core} compare both corrected
\grkan{} group counts against the MLP baseline on bpb and DCLM CORE.
The best corrected configuration is $g=4$: it reaches val\_bpb = 0.890 versus
0.8464 for MLP (gap $+0.044$, $\sim5.2\%$) and CORE = 0.128 versus 0.154
($\Delta=-0.026$).
Increasing the group count to $g=8$ \emph{hurts} both metrics
(val\_bpb = 0.9030, $+0.057$/$\sim6.7\%$; CORE = 0.116, $\Delta=-0.038$), so a
higher group count provides no benefit at this scale.
Against the stronger $g=4$ variant, MLP leads on 18 of the 22 CORE tasks
($g=4$ leads on \texttt{bigbench\_cs\_algorithms}, \texttt{agi\_eval\_lsat\_ar},
and \texttt{bigbench\_dyck\_languages}, with one tie); against $g=8$, MLP leads on
19 of 22.
The two consistent \grkan{} wins---\texttt{bigbench\_cs\_algorithms}
(0.432/0.434 vs.\ 0.417) and \texttt{agi\_eval\_lsat\_ar} (0.114/0.109 vs.\
0.011)---hold for both group counts, but the COPA advantage reported in the
preliminary version does not survive the denominator correction (COPA is now an
MLP win, 0.140 vs.\ 0.040/0.020).

\begin{table}[h]
\centering
\caption{%
  \grkan{} vs.\ MLP at GPT-2 scale: validation bpb
  (ClimbMix, $d_\text{model}=768$, 12 layers, 286M parameters, 2\,520 steps,
  H100 80GB HBM3, Muon optimizer).
  Canonical Safe Pad\'{e} rational formula; both \grkan{} rows are trained from
  scratch with the corrected kernel.
}
\label{tab:grkan_gpt2}
\begin{tabular}{lccc}
\toprule
Model & Val bpb $\downarrow$ & Gap vs.\ MLP & Tok/sec \\
\midrule
MLP transformer (SwiGLU) & \textbf{0.8464} & --- & ${\sim}508$K \\
\grkan{} ($g=4$, $P/Q=5/4$) & 0.890 & $+0.044$ ($+5.2\%$) & ${\sim}290$K \\
\grkan{} ($g=8$, $P/Q=5/4$) & 0.9030 & $+0.057$ ($+6.7\%$) & ${\sim}290$K \\
\bottomrule
\end{tabular}
\end{table}

\begin{table}[h]
\centering
\caption{%
  \grkan{} ($g=4$, $g=8$) vs.\ MLP at GPT-2 scale: DCLM CORE benchmark
  (22 ICL tasks, centered accuracy unless noted; higher is better).
  Selected per-task breakdown; full 22-task CSVs are released with the code.
  Canonical formula, corrected kernel.
}
\label{tab:grkan_core}
\begin{tabular}{lcccr}
\toprule
Task & \grkan{} ($g=4$) & \grkan{} ($g=8$) & MLP & Winner \\
\midrule
CORE (aggregate) & 0.128 & 0.116 & \textbf{0.154} & MLP \\
\midrule
HellaSwag (10-shot) & 0.103 & 0.094 & \textbf{0.140} & MLP \\
ARC-Easy & 0.369 & 0.335 & \textbf{0.422} & MLP \\
ARC-Challenge & 0.009 & $-0.009$ & \textbf{0.046} & MLP \\
PIQA & 0.313 & 0.312 & \textbf{0.357} & MLP \\
LAMBADA (accuracy) & 0.256 & 0.253 & \textbf{0.304} & MLP \\
COPA & 0.040 & 0.020 & \textbf{0.140} & MLP \\
BoolQ & $-0.081$ & $-0.167$ & \textbf{$-0.058$} & MLP \\
\midrule
bigbench\_cs\_algorithms & 0.432 & \textbf{0.434} & 0.417 & \grkan{} \\
agi\_eval\_lsat\_ar & \textbf{0.114} & 0.109 & 0.011 & \grkan{} \\
\bottomrule
\end{tabular}
\end{table}

\paragraph{Interpretation.}
The corrected \grkan{} trails MLP at every tested group count---by $+0.044$ bpb
($g=4$) to $+0.057$ bpb ($g=8$), and by $-0.026$ to $-0.038$ on CORE.
This is a larger gap than the $\sim1.2\%$ incorrectly reported in a preliminary
version, and the striking COPA advantage previously reported was an artifact of
the denominator bug.
The bpb gap is still smaller than the 26\% gap for \mlpedge{}, so rational
activations remain a less damaging direction than per-edge MLPs at this scale, but
the current result is negative relative to MLP.
Notably, the group count does \emph{not} behave as a free capacity knob: $g=8$ is
worse than $g=4$ on both bpb and CORE, and a $g=16$ run diverged to NaN, suggesting
the group expansion needs reoptimization rather than naive scaling at this depth.
Whether the gap closes with more training tokens, different initialization schemes
for the rational coefficients, or a re-tuned group count is unknown.
Each configuration here is a single run; we leave multi-seed replication, longer
training horizons, and hyperparameter exploration to future work.
% ── 6. Interpretability Analysis ──────────────────────────────────────────────
\section{Interpretability Analysis}
\label{sec:interpretability}

\subsection{Edge function reconstruction}

For each \kanlinear{} layer, we reconstruct all $d_\text{in}\!\times\!d_\text{out}$
learned 1D functions:
\begin{equation}
  \fij(x) = W_\text{base}[j,i]\cdot\text{SiLU}(x)
           + \sum_k W_\text{spline}[j,i,k]\cdot B_{i,k}(x)
\end{equation}
evaluating each on 200 points from its per-channel grid range.

\paragraph{MLP effective edge functions.}
For continuity with the earlier GuppyLM-style MLP comparison used in our
interpretability notebook, we extract an ``effective edge function''
$\fij^\text{MLP}$ from a 2-layer FFN:
\begin{equation}
  \fij^\text{MLP}(x) = \mathbf{W}_2[j,:]\cdot\text{ReLU}(\mathbf{W}_1[:,i]\cdot x + \mathbf{b}_1) + \mathbf{b}_2[j]
\end{equation}
with $\mathbf{W}_1\in\mathbb{R}^{768\times384}$ and $\mathbf{W}_2\in\mathbb{R}^{384\times768}$.
This is the scalar-valued contribution of input channel~$i$ to output channel~$j$
through the FFN, holding all other input channels at zero.
We sample the same 200-point range $[-2, 2]$ for each pair.

\subsection{Per-function metrics}

We compute three scalar metrics for every edge:

\textbf{Nonlinearity score (NLS):}
$\text{NLS} = \|f - f_\text{lin}\|_2 / (\|f\|_2 + \varepsilon)$,
where $f_\text{lin}$ is the best-fit affine function.
NLS~$\approx 0$ indicates a near-linear edge; NLS~$\approx 1$ indicates
strongly nonlinear behavior.

\textbf{Activity:}
$\text{Act} = \|f\|_2 / \sqrt{n}$.
Near-zero activity flags dead edges (prunable).

\textbf{Roughness:}
$\text{Rough} = \text{mean}(|\Delta^2 f|)$, the mean absolute second finite
difference.

\paragraph{Result summary.}
Across all 6 FFN layers and 884,736 edge functions:
\textbf{87.8\%} of edges have NLS $>0.1$ (meaningfully nonlinear),
and only \textbf{0.4\%} are effectively dead (activity $\leq 0.01$).
NLS is remarkably consistent across layers (median 0.279–0.282 per layer),
suggesting that nonlinearity is a stable property of KAN FFN layers rather
than concentrated in specific depths.
Figure~\ref{fig:nls} shows the NLS distributions for KAN vs.\ MLP
effective edge functions; a Mann-Whitney~U test gives $p < 10^{-300}$ for
KAN\,$>$\,MLP, consistent across all 6 layers individually.
This $p$-value is reported as a descriptive summary: edges within a layer share
the same B-spline grid and basis functions, so they are not mutually independent,
and the effective sample size is substantially smaller than 884,736.
The directional conclusion---KAN edges are quantitatively more nonlinear than MLP
edges at this scale---is supported by the per-layer consistency of the effect.
MLP edges are also predominantly nonlinear (median NLS\,=\,0.268 vs.\ 0.281 for KAN),
suggesting that the difference is quantitative rather than qualitative.

\subsection{Functional PCA}

We apply truncated SVD to the $(d_\text{out}\!\cdot d_\text{in}) \times n_\text{pts}$
curve matrix per layer (curves normalised to a common $[0,1]$ domain),
giving functional principal components analogous to fPCA.
Figure~\ref{fig:fpca} (Appendix) shows the top-4 components for all 6 layers.

\paragraph{Result.}
The top 4 functional principal components explain \textbf{99.9\%} of curve variance
in every layer (PC1: 60--62\%, PC2: 30--32\%, PC3: 7\%, PC4: $<$0.3\%).
We note that the finite B-spline basis bounds the edge-function space to a
small number of degrees of freedom, so finding 99.9\% variance in 4 components
is partly a consequence of the parametrization itself.
The substantive finding is \emph{which} shapes dominate: fPC1 is a smooth
monotone S-curve (sigmoid-like); fPC2--4 capture symmetric, asymmetric, and
higher-frequency perturbations.
The model converged to fewer functional archetypes than the available basis
functions allow, indicating that training collapses the accessible function
space to a small set of dominant motifs.

\subsection{MLPEdge preserves interpretability}

Figure~\ref{fig:edge_compare} plots $\fij(x)$ for the same $(i,j)$ pairs
from \kanprey{} and \mlpedge{}.
Both produce smooth, visualizable 1D curves; the \mlpedge{} curves are
comparable in smoothness but have global support (SiLU-based), unlike the
locally-supported B-spline functions.

\subsection{Closed-form approximation of edge functions}
\label{sec:symbolic}

A key interpretability promise of KANs is that learned edge functions can be
matched to closed-form expressions.
We test closed-form \emph{approximability} by fitting a six-element library to
the top-50 most active edges per layer: $\{$linear, quadratic, cubic, sigmoid,
tanh, $x\cdot\text{SiLU}(x)\}$.
For each edge, we report the best-fitting form by $R^2$ on 200 evaluation points
within the training activation range.
We do \emph{not} claim this constitutes symbolic regression in the full sense:
no search over expression trees is performed and the library is intentionally
small.

\paragraph{Result.}
\textbf{93.3\%} of the top-50 most active edges per layer admit a high-fidelity
fit with $R^2 > 0.99$ (mean $R^2 = 0.9961$, median $= 0.9993$).
The dominant form is cubic polynomial (91\%), with swish ($x \cdot
\text{SiLU}(x)$, 8\%) and sigmoid (1\%) accounting for the remainder.
Figure~\ref{fig:symbolic} (Appendix) shows representative edge functions with
their best-fit approximations.
These results confirm that the most active edge functions are well-approximated
by low-degree polynomials on the training domain, consistent with the KAN
interpretability motivation.
Crucially, however, high $R^2$ for a cubic fit is not discriminative: Taylor's
theorem guarantees that any smooth function on a bounded interval can be
approximated arbitrarily well by a sufficiently high polynomial, so finding that
cubic fits suffice does not identify the symbolic form.
Furthermore, the six-element library is far too small to constitute a genuine
symbolic search.
A rigorous test would require a grammar-based search (e.g., PySR~\citep{cranmer2023interpretable})
evaluated on held-out input ranges outside the training distribution.
We regard the present analysis as characterizing the \emph{complexity class} of
the learned functions (smooth, bounded, polynomial-like) rather than identifying
their symbolic forms.

\subsection{Domain-selective circuits}
\label{sec:circuits}

We register forward hooks on the layer-2 \kanlinear{} and record per-channel
activations over 100 fish-domain prompts and 100 generic prompts sampled from
the held-out test split.
For each edge $(i,j)$ we compute
$\Delta_{i,j} = f_{i,j}(\bar{x}^{\text{fish}}_i) - f_{i,j}(\bar{x}^{\text{gen}}_i)$,
where $\bar{x}^{\cdot}_i$ is the mean pre-KAN activation of input channel~$i$
averaged over all token positions and prompts in each domain.
Figure~\ref{fig:circuits} shows the resulting difference heatmap for layer~2.

\paragraph{Result.}
A sparse set of edges shows strong domain selectivity: the top-20 most
selective edges have $|\Delta| > 0.05$ while the global mean is $|\Delta| \approx 0.002$.
Input channel~147 has the highest mean $|\Delta|$ across all its output edges
(mean 0.016 vs.\ layer mean 0.005) and appears in 3 of the top 5 most
domain-selective pairs, identifying it as a domain-routing hub.
This edge-level circuit structure is a correlational observation enabled by the
KAN's explicit 1D decomposition --- an analysis that has no direct analogue in
MLP FFNs, where activations are entangled across input dimensions.

\subsection{Structured pruning}

The interpretability analysis enables principled pruning: edges with low
activity (small $\|f\|_2$) contribute negligibly to model output and can be
zeroed without retraining.
We evaluate two single-criterion pruning strategies:
(a)~\emph{activity-only}: prune edge $(i,j)$ if
    $\text{Act}_{i,j} < \tau_\text{act} \cdot \bar{\text{Act}}_\ell$
    where $\bar{\text{Act}}_\ell$ is the layer mean; and
(b)~\emph{NLS-only}: prune if $\text{NLS}_{i,j} < \tau_\text{nls}$.

\paragraph{Result.}
Activity-only pruning at $\tau_\text{act}=0.3$ removes \textbf{3.6\%} of edges
with only $\Delta\text{val\_loss} = +0.003$ zero-shot.
After 3{,}000 steps of fine-tuning, the pruned model recovers to
\textbf{0.297} val loss.
Note: this is compared against the \emph{uncorrected} \kanprey{} baseline
(0.402, 4\,000 steps, no masking); the corrected baseline is 0.2894, so the
pruned+finetuned model (0.297) does \emph{not} recover to baseline quality.
This correctly frames the result: 3.6\% sparsity incurs a permanent
$+0.008$ val loss cost after fine-tuning.
NLS-only pruning tells the complementary story: removing edges with
NLS~$<0.05$ (just 2.8\% of edges) causes an immediate $+0.042$ degradation,
confirming that even the least-nonlinear KAN edges carry meaningful signal.
Figure~\ref{fig:pruning} (Appendix) plots the full quality--sparsity curves.

% ── 7. Discussion ─────────────────────────────────────────────────────────────
\section{Discussion}
\label{sec:discussion}

\subsection{Why MLPEdge works}

The KAN topology---additive decomposition of each output feature over all
input channels via independent 1D functions---remains the key architectural
inductive bias inside the KAN-family design space.
At small scale (384 hidden units, 60\,K training samples), a tiny MLP per
edge improves over the tested B-spline KAN variants even without a grid update
mechanism, suggesting that the B-spline basis itself is not the only route to
useful edge-wise functions.

\subsection{Honest tradeoffs of MLPEdge}

\textbf{Advantages.}
No grid mechanism (no extra hyperparameters, no update step), and within the
tested KAN-family variants, better validation loss and stronger heuristic
prompt behavior than the B-spline baselines in Table~\ref{tab:main}.

\textbf{Disadvantages.}
B-spline knots provide locally-supported basis functions, giving the spline
path an inductive bias for smooth, locally-adapting functions;
\mlpedge{} uses SiLU, which has global support.
Symbolic regression of learned edge functions is harder: B-splines can be
approximated by simple polynomials per segment, while SiLU-based MLPs have
no closed-form simplification.
The connection to the Kolmogorov-Arnold theorem is equally valid: a SiLU-based
MLP mapping $\mathbb{R}\to\mathbb{R}$ is a continuous univariate function and
satisfies the theorem's requirements.
The practical distinction is that B-spline edges have a known, symbolic basis
suited for regression, while SiLU-MLP edges have no closed-form simplification.

\textbf{Scale caveat.}
The small-scale improvement of \mlpedge{} over the tested B-spline KAN
variants does not close the gap to the standard MLP baseline, and at GPT-2
size (Section~\ref{sec:scaling}) the gap widens substantially: the MLP
baseline outperforms the parameter-matched \mlpedge{} by 4.4 perplexity
points while also running 11\% faster.
We therefore do not claim that \mlpedge{} generalises beyond the
\guppylm{} testbed.

\subsection{GR-KAN at GPT-2 Scale}

The corrected \grkan{} scale comparison in Section~\ref{sec:grkan_gpt2} is
cautionary.
Under the canonical Safe Pad\'{e} formula and a group-count sweep, the best
\grkan{} configuration ($g=4$) reaches val\_bpb = 0.890 versus 0.8464 for the
MLP baseline at 286M parameters on ClimbMix---a 5.2\% gap---while $g=8$ is worse
(0.9030, 6.7\%).
On DCLM CORE (22 ICL tasks), $g=4$ scores 0.128 and $g=8$ scores 0.116 versus
0.154 for MLP ($\Delta=-0.026$ and $-0.038$), with MLP leading on 18 of 22 tasks
against $g=4$ and no COPA advantage at either group count.
The bpb gap is smaller than the 26\% \mlpedge{} gap at GPT-2-small scale, so
rational activations remain a plausible but unproven route for future work.
These results do not establish that \grkan{} is competitive with MLP, nor that
group-shared Pad\'{e} activations scale cleanly.

Because all smaller pre-fix \grkan{} runs used the wrong denominator formula,
we do not draw conclusions from their reported validation loss, modular-staging
behavior, or external comparisons.
Those rows require fresh denominator-corrected reruns before they can re-enter
the paper.

\subsection{Relation to PowerMLP}

\citet{tang2025powermlp} show algebraic equivalence between B-splines and
ReLU-power compositions, but do not train or evaluate per-edge MLP functions
in a language model.
\mlpedge{} is independently motivated and evaluated here;
the equivalence result provides a theoretical connection but does not
anticipate our empirical findings.

% ── 8. Limitations ────────────────────────────────────────────────────────────
\section{Limitations}
\label{sec:limitations}

\guppylm{} is a controlled testbed, not a production-scale model.
Our GPT-2 scale experiment (Section~\ref{sec:scaling}) shows that
\mlpedge{}'s advantage does \emph{not} transfer to $d_\text{model}=768$:
the per-edge factored structure falls 26\% behind the MLP baseline in
perplexity at matched parameter count.
Whether a larger per-edge hidden size or different initialisation would
recover competitive performance for \mlpedge{} at scale remains open.
Section~\ref{sec:grkan_gpt2} shows that canonical \grkan{} with rational
activations also trails MLP at 286M parameters at every tested group count:
$+0.044$ ($g=4$) to $+0.057$ ($g=8$) bpb, $-0.026$ to $-0.038$ DCLM CORE,
and lower throughput.
Although this bpb gap is smaller than the \mlpedge{} gap, it is still a negative
result relative to the MLP baseline and cannot support a \grkan{} architecture
claim.
The current corrected comparison is a single run per group count ($g \in \{4,8\}$); additional seeds, other
initialization schemes, and longer training horizons remain open.
The keyword-match evaluation metric on \guppylm{} is coarse; the GPT-2
scale experiment uses bpb (ClimbMix) and DCLM CORE (ICL benchmark).
Interpretability results apply to B-spline KANLinear; \mlpedge{} functions
are visualizable but harder to analyze symbolically.
The domain-selective circuit analysis (Section~\ref{sec:circuits}) is
correlational.
We tested causal relevance by scaling the spline output of the top-5 most
domain-selective edges by $\alpha \in \{0, 0.5, 2, 5\}$ (zeroing at
$\alpha=0$, amplifying at $\alpha=5$) and measuring the resulting change in
validation loss on 100 fish-domain vs.\ 100 generic prompts.
The maximum observed asymmetry (|$\Delta\text{loss}_\text{fish}
- \Delta\text{loss}_\text{generic}|$) across all $\alpha$ values is
$\mathbf{<0.004}$ --- comparable to noise and far below what is needed to
claim selective behavioral impact.
This null result likely reflects redundancy across the 384 input channels and
6 transformer layers; single-edge interventions are insufficient.
Establishing causal circuit roles would require larger multi-edge ablations or
activation patching across the full residual stream.
Training logs for early experiments were not saved step-by-step;
Figure~\ref{fig:loss_curves} uses logged val\_loss from re-runs with the
corrected training loop.

% ── 9. Related Work ───────────────────────────────────────────────────────────
\section{Related Work}
\label{sec:related}

\paragraph{KAN and variants.}
\citet{liu2024kan} introduce KANs and demonstrate symbolic regression on
scientific functions. \citet{yang2024kat} extend KANs to attention layers.
\citet{liu2024kan2} (KAN 2.0) add MultKAN and grid extension techniques.
\citet{bodner2024convolutional} and others apply KANs to vision tasks.

\paragraph{Efficient KAN.}
\citet{blealtan2024efficientkan} implement \kanlinear{} efficiently using
pre-computed B-spline grids; we build on this implementation.

\paragraph{Additive models and interpretability.}
Neural additive models~\citep{agarwal2021neural} represent each feature by a
learned shape function, sharing the ``one function per input dimension''
interpretability goal.
KANs extend this to the edge level, yielding a richer family of decompositions.

\paragraph{PowerMLP.}
\citet{tang2025powermlp} prove algebraic equivalence between KAN spline
edges and ReLU-power compositions, providing theoretical motivation
for replacing splines with polynomial networks.
\mlpedge{} pursues a similar substitution but uses learned per-edge MLPs and
evaluates the approach in a language model setting.

% ── 10. Conclusion ────────────────────────────────────────────────────────────
\section{Conclusion}
\label{sec:conclusion}

We presented the first empirical interpretability study of KAN edge functions
in a language model, finding that 87.8\% of 884,736 edges are meaningfully
nonlinear (NLS\,$>0.1$) and only 0.4\% are dead.
At \guppylm{} scale, the function-basis screen finds that a grouped Chebyshev
degree-3 FFN reaches $0.2770_{\pm 0.0003}$ validation loss over three seeds,
matching the same-protocol MLP mean ($0.2770_{\pm 0.0008}$) but running slower
in the seed-42 inference sweep (550\,ms vs.\ 309\,ms).
Corrected canonical \grkan{} remains close ($0.2780_{\pm 0.0005}$), and a
square-denominator rational ablation is similar ($0.2780_{\pm 0.0009}$), but no
tested KAN-family FFN dominates MLP on both quality and latency.
\mlpedge{} remains the strongest visualizable edge-function replacement among
the B-spline and \mlpedge{} rows: it reaches $0.2854_{\pm 0.0005}$ validation
loss at $h=5$, compared with $0.2883_{\pm 0.0005}$ for the standard-attention
B-spline baseline and $0.2867_{\pm 0.0006}$ for \katprey{}.
At GPT-2 scale (124M parameters, Wikitext-103), the parameter-matched
\mlpedge{} ($h=8$) trails the MLP baseline by a substantial margin
(perplexity 20.96 vs.\ 16.58), indicating that the per-edge factored
structure does not scale straightforwardly beyond the small-model regime
explored here.

All pre-fix \grkan{} runs outside the corrected \guppylm{} rerun are excluded
from the evidence base because they used the wrong rational denominator formula.
The removed modular-staging rows still require fresh denominator-corrected
reruns; the old d12 group-count rows, however, have been replaced by the
corrected from-scratch runs reported next.
We additionally scale \grkan{} to 286M parameters on ClimbMix
(Section~\ref{sec:grkan_gpt2}), finding that the best corrected configuration
($g=4$) reaches val\_bpb = 0.890 versus 0.8464 for the MLP baseline---a 5.2\%
gap, widening to 6.7\% at $g=8$---after
engineering a fused Triton backward kernel for the group-rational activation
(${\sim}3\times$ faster than the unfused rational baseline; ${\sim}290$K
tok/s versus ${\sim}508$K tok/s for MLP).
On DCLM CORE (22 ICL tasks), $g=4$ scores 0.128 and $g=8$ scores 0.116 versus
0.154 for MLP ($\Delta=-0.026$ and $-0.038$), with MLP leading on 18 of 22 tasks
and no COPA advantage
under the corrected denominator.
Thus, the corrected \grkan{} result is a useful negative/cautionary scale
result: rational activations appear less damaging than the per-edge \mlpedge{}
factorization at this scale, but they do not yet match the MLP baseline.

The \guppylm{} testbed and all model code are publicly available:
\url{https://github.com/HCAI-USP/kan-guppylm}.

% ── Acknowledgments ───────────────────────────────────────────────────────────
\begin{ack}
This work was carried out at the Human-Centered Artificial Intelligence group
(HCAI) of the University of São Paulo.
We thank the GuppyLM dataset contributors and the EfficientKAN open-source
implementation.
Small-scale experiments were run on a personal Apple Silicon workstation.
The Wikitext-103 GPT-2 scale experiments used an NVIDIA RTX~A6000 GPU via RunPod community cloud; the ClimbMix \grkan{} scale comparison used an H100 80GB HBM3 GPU.
No external compute grants were used.
\end{ack}

% ── References ────────────────────────────────────────────────────────────────
\bibliographystyle{plainnat}
\bibliography{references}

% ── Appendix ──────────────────────────────────────────────────────────────────
\appendix

\section{Architecture Diagram}
\label{app:architecture}

\begin{figure}[h]
\centering
% Architecture schematic generated from current model definitions.
\includegraphics[width=\linewidth]{figures/architecture.pdf}
\caption{%
  \textbf{Three FFN variants.}
  \emph{Left}: Standard MLP FFN (two linear layers + ReLU).
  \emph{Center}: KANLinear --- each edge $\fij$ is a B-spline function with
  learnable coefficients and adaptive knots.
  \emph{Right}: \mlpedge{} --- each edge $\fij$ is a tiny MLP
  ($\mathbb{R}\to\mathbb{R}$), realized via the same einsum structure.
}
\label{fig:architecture}
\end{figure}

\section{Training Loss Curves}
\label{app:loss}

\begin{figure}[h]
\centering
% Loss curves generated from current training logs.
\includegraphics[width=0.7\linewidth]{figures/loss_curves.pdf}
\caption{%
  Validation loss vs.\ training step for all four models.
  The prompt-masking fix (dashed vertical line at step 0 of KAT v2 and
  \mlpedge{} runs) accounts for the discontinuous improvement relative to
  the initial \kanprey{} run.
}
\label{fig:loss_curves}
\end{figure}

\section{Spline vs.\ MLPEdge Edge Visualization}
\label{app:edge_viz}

\begin{figure}[h]
\centering
% Edge comparison figure generated from current interpretability artifacts.
\includegraphics[width=0.8\linewidth]{figures/edge_compare.pdf}
\caption{%
  Learned 1D edge functions $\fij(x)$ for the same $(i, j)$ pairs from
  layer~0 of \kanprey{} (B-spline, \emph{left}) and \mlpedge{}
  (\emph{right}).
  Both produce smooth, visualizable curves; \mlpedge{} functions have global
  support while B-spline functions have locally-supported knots (vertical
  tick marks, bottom).
}
\label{fig:edge_compare}
\end{figure}

\section{MLPEdge Hidden-Size Ablation}
\label{app:ablation}

\begin{figure}[h]
\centering
% Ablation figure generated from current MLPEdge sweep artifacts.
\includegraphics[width=0.65\linewidth]{figures/ablation.pdf}
\caption{%
  \mlpedge{} hidden-size ablation.
  Left axis: validation loss; right axis: inference latency.
  $h=1$ achieves the best validation loss (0.2828); $h=2$ is the practical
  sweet spot (0.2836, 354\,ms).
  We use $h=5$ as our main comparison point for parameter-count matching
  with $n_\text{basis}=5$ B-spline models.
}
\label{fig:ablation}
\end{figure}

\section{NLS Distribution and MLP Comparison}
\label{app:nls}

\begin{figure}[h]
\centering
\includegraphics[width=0.7\linewidth]{figures/nls_hist.pdf}
\caption{%
  Overlaid histograms of nonlinearity score (NLS) for KAN edge functions
  (884,736 edges across 6 FFN layers, blue) and MLP effective edge functions
  (orange). Log-scale y-axis.
  Mann-Whitney~U test: $p < 10^{-300}$ (KAN\,$>$\,MLP).
  Both distributions peak well above 0.1, with KAN edges
  slightly more nonlinear (median 0.281 vs.\ 0.268).
}
\label{fig:nls}
\end{figure}

\section{Functional PCA Components}
\label{app:fpca}

\begin{figure}[ht]
\centering
\includegraphics[width=\linewidth]{figures/fpca_all_layers.pdf}
\caption{%
  Top-4 functional principal components for all 6 KAN FFN layers of \kanprey{}.
  PC1--4 together explain \textbf{99.9\%} of curve variance in every layer
  (PC1 alone: 60--62\%).
  Note: the finite B-spline basis bounds the edge-function space to a small
  number of degrees of freedom; the 4-PC result reflects that training
  further collapses this space to dominant archetypes (see Section~\ref{sec:interpretability}).
}
\label{fig:fpca}
\end{figure}

\section{Domain-Selective Circuit Heatmap}
\label{app:circuits}

\begin{figure}[ht]
\centering
\includegraphics[width=\linewidth]{figures/circuits_heatmap.pdf}
\caption{%
  Per-edge output difference $\Delta_{i,j} = f_{i,j}(\bar{x}^\text{fish}_i)
  - f_{i,j}(\bar{x}^\text{gen}_i)$ for layer~2 of \kanprey{}.
  \emph{Left}: mean edge outputs for fish-domain prompts.
  \emph{Centre}: mean edge outputs for generic prompts.
  \emph{Right}: difference (fish\,$-$\,generic). Red = fish-selective,
  blue = generic-selective.
  A sparse set of edges ($|\Delta| > 0.05$) shows strong domain selectivity;
  input channel 147 appears in 3 of the 5 most domain-selective pairs.
}
\label{fig:circuits}
\end{figure}

\section{Structured Pruning Quality--Sparsity Curves}
\label{app:pruning}

\begin{figure}[ht]
\centering
\includegraphics[width=\linewidth]{figures/pruning_sparsity.pdf}
\caption{%
  Zero-shot val loss vs.\ fraction of edges pruned for activity-only
  (\emph{left}) and NLS-only (\emph{right}) pruning strategies.
  Activity-only pruning removes 3.6\% of edges at $+0.003$ val loss;
  NLS-only pruning causes steep degradation even at low sparsity,
  confirming that nonlinear edges are functionally load-bearing.
  Dashed line: baseline; dotted line: baseline $+0.01$.
}
\label{fig:pruning}
\end{figure}


\section{Closed-Form Approximation of Edge Functions}
\label{app:symbolic}

\begin{figure}[ht]
\centering
\includegraphics[width=\linewidth]{figures/symbolic_regression.pdf}
\caption{%
  Representative KAN edge functions (blue) from layer~0 of \kanprey{} with
  their best-fit closed-form approximations (orange dashed).
  Candidate forms: linear, quadratic, cubic, sigmoid, tanh, $x\cdot\text{SiLU}(x)$.
  $R^2$ is reported for each fit within the training activation range.
  93.3\% of the top-50 most active edges per layer achieve $R^2 > 0.99$
  (mean $R^2 = 0.9961$); the dominant form is cubic polynomial (91\%).
  Note: high $R^2$ for a cubic on a bounded interval reflects the smoothness
  of the learned functions, not a definitive symbolic identification
  (see Section~\ref{sec:symbolic}).
}
\label{fig:symbolic}
\end{figure}


\section{Reproducibility Checklist}
\label{app:reprod}

This section documents the exact software, hardware, datasets, commands, and seeds needed to reproduce the main experiments.

\subsection{Code and Revisions}

\begin{itemize}
  \item \textbf{kan-guppylm} (GuppyLM runs, MLPEdge GPT-2 Wikitext-103):
    \url{https://github.com/HCAI-USP/kan-guppylm},
    commit \texttt{3c8ea92} (2026-06-01).
  \item \textbf{nanokan} (GR-KAN ClimbMix d12 runs):
    \url{https://github.com/ACS-USP/nanokan},
    commits \texttt{8b21df2} ($g=4$ run) and \texttt{0e0bea5} ($g=8$ run); the
    canonical Safe Pad\'{e} denominator is provided by the pinned
    \texttt{rational\_kat\_cu} kernel below.
  \item \textbf{nanochat} (training harness dependency):
    \url{https://github.com/karpathy/nanochat},
    commit \texttt{dc54a1a} (2026-05-22).
  \item \textbf{rational\_kat\_cu} (rational activation kernels):
    \url{https://github.com/Adamdad/rational\_kat\_cu},
    commit \texttt{41a20b5}, with the corrected Safe Pad\'{e} denominator and fused Triton backward kernels.
\end{itemize}

\subsection{Rational Denominator Correction}

All pre-correction GR-KAN d12 checkpoints are excluded from the publication evidence
except as historical invalidation context.
The bug and fix are as follows:
\begin{itemize}
  \item \textbf{Canonical (correct):} $Q(x) = 1 + |b_0 x + b_1 x^2 + b_2 x^3 + b_3 x^4|$.
  \item \textbf{Retracted (bugged):} $Q(x) = 1 + |b_0||x| + |b_1||x|^2 + |b_2||x|^3 + |b_3||x|^4$.
\end{itemize}
Gradient consistency of the corrected kernel was verified via \texttt{torch.autograd.gradcheck}
on H100 hardware (max errors: $\mathrm{d}x = 3.81\times10^{-6}$, $\mathrm{d}a = 8.58\times10^{-6}$,
$\mathrm{d}b = 9.54\times10^{-6}$, float32).

\subsection{Environment Versions}

\begin{table}[h]
\centering
\caption{Software environment for each experiment family.}
\label{tab:env}
\begin{tabular}{lccc}
\toprule
  & GuppyLM Track A & MLPEdge GPT-2 & GR-KAN ClimbMix d12 \\
\midrule
OS & Darwin 25.4.0 arm64 & Ubuntu 22.04 & Ubuntu 22.04 \\
Python & 3.13.5 & 3.10.12 & 3.10.12 \\
PyTorch & 2.12.0 (MPS) & 2.4.1 (CUDA 12.4) & 2.5.1 (CUDA 12.4) \\
Triton & --- & --- & 3.1.0 \\
Flash Attn & --- & --- & FA3 (Hopper SM90) \\
Muon & --- & --- & Newton-Schulz-5 \\
\bottomrule
\end{tabular}
\end{table}

\subsection{Hardware}

\begin{itemize}
  \item \textbf{GuppyLM Track A}: Apple M4 Pro, 14 logical cores, 25.77~GB unified memory, MPS backend.
    Training wall-clock: $\sim$10\,min per 8,000-step run.
  \item \textbf{MLPEdge GPT-2 Wikitext-103}: single NVIDIA RTX~A6000 48~GB (RunPod community cloud).
    Training wall-clock: 13.3\,h (MLP), 14.9\,h (\mlpedge{}).
  \item \textbf{GR-KAN ClimbMix d12}: single NVIDIA H100 80~GB HBM3 (RunPod community cloud).
    Training wall-clock: $\sim$2.0\,h per run (2,520 steps); step time
    $\sim$1.76\,s, throughput $\sim$290K\,tok/s, $\sim$88\% BF16 MFU
    (steady state, from the run logs).
    Eval hardware: NVIDIA RTX PRO 4500 Blackwell 32~GB (separate pod).
\end{itemize}

\subsection{Dataset Provenance and Tokenizers}

\begin{itemize}
  \item \textbf{GuppyLM}: \texttt{arman-bd/guppylm-60k-generic} (HuggingFace;
    MIT license, DOI \texttt{10.57967/hf/8339}).
    Tokenizer: \texttt{tokenizer.json} (GPT-2 BPE, vocab 2,393).
    Max sequence length: 128 tokens.
    Objective: assistant-token-only loss; prompt token targets set to $-100$.
  \item \textbf{MLPEdge GPT-2 Wikitext-103}: \texttt{Salesforce/wikitext},
    configuration \texttt{wikitext-103-v1} (Merity et al., 2017;
    CC BY-SA 3.0 / GFDL).
    GPT-2 BPE tokenizer (vocab 50,257).
  \item \textbf{GR-KAN ClimbMix d12}: nanochat's ClimbMix corpus, traced by
    nanochat to NVIDIA Nemotron-ClimbMix (400B tokens; CC BY-NC 4.0;
    research and development use).
    RustBPE tokenizer trained on 30 ClimbMix shards,
    \texttt{--max-chars=2B --vocab-size=32768}.
    Tokenizer artifacts: \texttt{tokenizer.pkl} + \texttt{token\_bytes.pt}.
    \textbf{Evaluation guard}: the eval pod uses the exact same tokenizer artifact
    from training (30-shard build), not a reduced-shard rebuild.
\end{itemize}

\subsection{GuppyLM Track A Commands and Seeds}

All cleanup runs: 8,000 steps, batch size 32, \texttt{kanprey.train} defaults
(AdamW, cosine LR, 200 warmup steps). Seed list: $\{42, 43, 44\}$.

\begin{verbatim}
# Same-protocol MLP baseline
uv run python -m kanprey.train --model mlp --steps 8000 \
  --batch-size 32 --checkpoint-dir checkpoints/mlp_s42 --seed 42
# ... repeated for seeds 43, 44

# B-spline KAN grid=2
uv run python -m kanprey.train --model kan --grid-size 2 --steps 8000 \
  --batch-size 32 --checkpoint-dir checkpoints/kan_grid2_s42 --seed 42
# ... repeated for seeds 43, 44

# Parameter-matched MLPEdge h=8
uv run python -m kanprey.train --model mlpedge --mlp-edge-hidden 8 \
  --steps 8000 --batch-size 32 \
  --checkpoint-dir checkpoints/mlpedge_h8_s42 --seed 42
# ... repeated for seeds 43, 44

# Corrected canonical GR-KAN
uv run python -m kanprey.train --model grkan --steps 8000 \
  --batch-size 32 --checkpoint-dir checkpoints/grkan_corrected_s42 \
  --seed 42
# ... repeated for seeds 43, 44
\end{verbatim}

Existing KAT and MLPEdge $h=5$ rows use the same seeds from \texttt{checkpoints/kat\_s*} and
\texttt{checkpoints/mlpedge\_s*}.

\subsection{GuppyLM Function-Basis Screen}

The function-basis screen first ran the Gaussian/RBF microbenchmark gate:
\begin{verbatim}
uv run python scripts/benchmark_basis_kernels.py \
  --shape 4096x384 --device auto --dtype float32 \
  --centers 8 --width-scale 1.5 --repeats 20 \
  --output results/basis_gaussian_microbench_4096x384_h15.json
\end{verbatim}

The Stage-1 pilot matrix used 2\,000 steps, batch size 32, seed 42, and
\texttt{checkpoints/basis\_pilots/<name>} output directories.
The promoted Stage-2 confirmation rows used:
\begin{verbatim}
# Square-denominator GR-KAN confirmation
uv run python -m kanprey.train --model grkan --steps 8000 \
  --batch-size 32 --seed 42 --grkan-denominator square \
  --grkan-m 5 --grkan-n 4 --grkan-groups 8 \
  --checkpoint-dir checkpoints/basis_confirm/grkan_square_p5q4_g8_s42
# ... repeated for seeds 43, 44

# Chebyshev degree-3 grouped basis confirmation
uv run python -m kanprey.train --model basis --steps 8000 \
  --batch-size 32 --seed 42 --basis-family chebyshev \
  --basis-degree 3 --basis-groups 8 --basis-input-norm tanh \
  --checkpoint-dir checkpoints/basis_confirm/cheb_d3_g8_s42
# ... repeated for seeds 43, 44
\end{verbatim}

Primary artifacts are \texttt{results/basis\_stage1\_summary.json},
\texttt{results/basis\_stage2\_summary.json},
\texttt{results/basis\_stage2\_run\_manifest.json},
\texttt{results/basis\_stage2\_eval.json}, and
\texttt{results/basis\_stage2\_diagnostics.json}.

\subsection{GuppyLM Unified Evaluation}

\begin{verbatim}
uv run python eval_cases.py --json --max-tokens 64 \
  --run MLP=checkpoints/mlp_s42/best.pt \
  --run KAN=checkpoints/kan_grid2_s42/best.pt \
  --run KAT=checkpoints/kat_s42/best.pt \
  --run MLPEdgeH5=checkpoints/mlpedge_s42/best.pt \
  --run MLPEdgeH8=checkpoints/mlpedge_h8_s42/best.pt \
  --run GRKAN=checkpoints/grkan_corrected_s42/best.pt \
  --run GRKANSquare=checkpoints/basis_confirm/grkan_square_p5q4_g8_s42/best.pt \
  --run ChebD3=checkpoints/basis_confirm/cheb_d3_g8_s42/best.pt
\end{verbatim}

Settings: device \texttt{mps}, tokenizer \texttt{tokenizer.json},
temperature 0.7, top-$k$ 50, max new tokens 64, 16 prompts.
The 16-prompt score is a pass/fail heuristic on manual criteria;
it is not a benchmark and should be interpreted with binomial uncertainty.
Latency: mean and median across 16 prompts $\times$ 64 generated tokens.

\subsection{MLPEdge GPT-2 Wikitext-103}

\begin{itemize}
  \item Model config: $d_\text{model}=768$, 12 heads, 12 layers, context 1,024.
  \item MLP FFN: $\text{Linear}(768{\to}3072) \to \text{GELU} \to \text{Linear}(3072{\to}768)$ (4$\times$ expansion).
  \item \mlpedge{} $h=8$: per-edge MLP hidden size 8, parameter-matched to MLP.
  \item Training: 20,000 steps, effective batch 128 (gradient accumulation $4 \times 32$),
    cosine LR $3\times10^{-4} \to 3\times10^{-5}$, 2,000 warmup, AdamW,
    gradient checkpointing enabled.
  \item Reported: best validation perplexity across all eval checkpoints.
\end{itemize}

\subsection{GR-KAN ClimbMix d12}

\begin{itemize}
  \item Harness: \texttt{nanokan} wrapping nanochat with \texttt{--ffn-type} switch.
  \item Model tags: \texttt{d12-mlp}, \texttt{d12-grkan-g4}, \texttt{d12-grkan-g8}.
  \item Config: $d_\text{model}=768$, 6 heads, 12 layers, sequence length 2,048,
    $\sim$286M parameters.
  \item GR-KAN: group counts $g \in \{4, 8\}$ (best: $g=4$; $g=16$ diverged), $P=5$, $Q=4$, canonical Safe Pad\'{e} denominator.
    Rational init: \texttt{identity} (rat1), \texttt{swish} (rat2).
  \item MLP baseline: $\text{Linear}(768{\to}3072) \to \text{SwiGLU} \to \text{Linear}(3072{\to}768)$.
  \item Training: 2,520 steps, effective batch 256 seq $\times$ 2,048 tok = 524,288 tok/step,
    $\sim$1.32B total tokens. Optimizer: Muon (Newton-Schulz-5)
    for 2D weights, AdamW for scalars and rational coefficients.
    BF16 mixed-precision, Flash Attention 3.
  \item Eval: BPB and DCLM CORE (22 ICL tasks, centered accuracy).
    Full per-task CSV available at \url{https://github.com/ACS-USP/nanokan}.
  \item RunPod commands:
\begin{verbatim}
python3 scripts/runpod_launch.py test
python3 scripts/runpod_launch.py train --model grkan \
  --volume-id <volume_id> --steps 2520
python3 scripts/runpod_launch.py eval \
  --model-tags d12-grkan-g4 d12-grkan-g8 d12-mlp --secure
\end{verbatim}
\end{itemize}

\subsection{Checkpoint Lineage}

\begin{table}[h]
\centering
\small
\caption{Checkpoint paths for every result row.}
\label{tab:checkpoints}
\begin{tabular}{@{}lp{0.58\linewidth}@{}}
\toprule
Result row & Checkpoints \\
\midrule
GuppyLM MLP & \texttt{mlp\_s42}, \texttt{mlp\_s43}, \texttt{mlp\_s44} \\
GuppyLM KAN grid=2 & \texttt{kan\_grid2\_s42}, \texttt{kan\_grid2\_s43}, \texttt{kan\_grid2\_s44} \\
GuppyLM KAT grid2+KAT & \texttt{kat\_s42}, \texttt{kat\_s43}, \texttt{kat\_s44} \\
GuppyLM MLPEdge $h=5$ & \texttt{mlpedge\_s42}, \texttt{mlpedge\_s43}, \texttt{mlpedge\_s44} \\
GuppyLM MLPEdge $h=8$ & \texttt{mlpedge\_h8\_s42}, \texttt{mlpedge\_h8\_s43}, \texttt{mlpedge\_h8\_s44} \\
GuppyLM corrected GR-KAN & \texttt{grkan\_corrected\_s42}, \texttt{grkan\_corrected\_s43}, \texttt{grkan\_corrected\_s44} \\
GuppyLM \grkan{} square denom. & \texttt{basis\_confirm/grkan\_square\_p5q4\_g8\_s42}, \texttt{basis\_confirm/grkan\_square\_p5q4\_g8\_s43}, \texttt{basis\_confirm/grkan\_square\_p5q4\_g8\_s44} \\
GuppyLM \basiskan{} Cheb $d=3,g=8$ & \texttt{basis\_confirm/cheb\_d3\_g8\_s42}, \texttt{basis\_confirm/cheb\_d3\_g8\_s43}, \texttt{basis\_confirm/cheb\_d3\_g8\_s44} \\
MLPEdge GPT-2 MLP & \texttt{checkpoints/scale/mlp/best.pt} \\
MLPEdge GPT-2 $h=8$ & \texttt{checkpoints/scale/mlpedge\_h8/best.pt} \\
GR-KAN d12 ($g=4$) & \nolinkurl{checkpoints/nanochat/b-g4-20260610T193454Z-full-d12-grkan-g4/d12-grkan-g4/model_002520.pt} plus \texttt{meta\_002520.json}, \texttt{tokenizer.pkl}, and \texttt{token\_bytes.pt} \\
GR-KAN d12 ($g=8$) & \nolinkurl{checkpoints/nanochat/20260608T193716Z-d12-grkan-g8-corrected-r2/d12-grkan-g8-corrected-r2/model_002520.pt} plus matching metadata and tokenizer \\
d12 MLP & \texttt{checkpoints/nanochat/d12-mlp/model\_002520.pt} plus matching tokenizer and eval artifacts \\
\bottomrule
\end{tabular}
\end{table}

\subsection{Statistical Reporting}

\begin{itemize}
  \item GuppyLM validation loss: mean $\pm$ std over 3 seeds for each model row.
  \item GuppyLM 16-prompt heuristic: pass/fail score; not a benchmark.
    Binomial standard error ${\sim}\pm2.5$ prompts at $n=16$.
    Differences of 1--2 correct responses are not statistically meaningful.
  \item MLPEdge GPT-2 Wikitext-103: single-run comparison (no multi-seed replication yet).
  \item GR-KAN ClimbMix d12: single corrected run per group count ($g=4$, $g=8$).
    DCLM CORE per-task results are reported with the aggregate; per-task
    advantage claims would require confidence intervals and predeclared analysis.
\end{itemize}

\subsection{Known Limitations Preserved in the Manuscript}

\begin{enumerate}
  \item No tested KAN-family GuppyLM model dominates MLP on both validation loss and latency; Chebyshev $d=3,g=8$ matches the MLP validation-loss mean but is slower.
  \item Corrected GuppyLM GR-KAN is close to MLP but not a replacement claim; the square-denominator rational ablation is also close but not a quality/latency win.
  \item Parameter-matched MLPEdge $h=8$ does not close the MLP gap at GuppyLM scale.
  \item MLPEdge $h=8$ fails to scale cleanly to GPT-2-small on Wikitext-103 (26\% perplexity gap).
  \item Corrected canonical GR-KAN trails MLP at 286M parameters on BPB, DCLM CORE, and throughput at every tested group count ($g=4$ best, $g=8$ worse, $g=16$ diverged).
  \item Pre-fix GR-KAN checkpoints outside the corrected GuppyLM rerun remain excluded from the evidence base until corrected reruns replace them.
  \item Domain-selective edge analysis is correlational; causal circuit claims are not supported.
\end{enumerate}
\end{document}
